% essaymaster paper skeleton — arXiv-style preprint (self-contained arxiv.sty).
% Fill every <PLACEHOLDER>; delete guidance comments as sections are written.
% Two-column densification (optional): uncomment the fontsize/geometry block.
\documentclass{article}          % add [twocolumn] for the dense variant
\usepackage{arxiv}
% \usepackage[fontsize=9.4pt]{fontsize}
% \geometry{a4paper, left=0.65in, right=0.65in, top=0.8in, bottom=0.8in}
% \setlength{\columnsep}{0.24in}
\usepackage[utf8]{inputenc}
\usepackage[T1]{fontenc}
\usepackage{hyperref}
\usepackage{url}
\usepackage{booktabs}
\usepackage{amsmath}
\usepackage{amssymb}
\usepackage{graphicx}
\usepackage{natbib}
\usepackage{xcolor}
\usepackage{microtype}

% Figures ship as PDF (preferred) with PNG fallback; build.sh converts the SVGs.
\graphicspath{{figures/}}
\hypersetup{colorlinks=true, linkcolor=blue!55!black, citecolor=blue!45!black, urlcolor=blue!55!black}
\bibliographystyle{plainnat}

\title{<TITLE: system name + what it is — no adjective pile-ups>}
\author{<AUTHOR>\\ \texttt{<email>}}
\date{\today}
\preprintlabel{<Short running label>: Preprint}

% Every unsourced number in the text uses \todo{} — loud, red, greppable.
% The pre-submission gate is: zero \todo{} left, or each remaining one is
% explicitly framed as future work in Limitations.
\newcommand{\todo}[1]{\textcolor{red}{[\textbf{TODO:} #1]}}

\begin{document}
\maketitle

\begin{abstract}
% Written LAST. Structure: problem + why it matters (1-2 sentences); what we built
% (1-2); headline measured results with qualifiers IN the sentence; the takeaway.
% Unrun measurements that matter get declared here, not hidden.
<ABSTRACT>
\end{abstract}

\section{Introduction}
\label{sec:intro}
% The case for the work; why prior approaches don't fit; then the explicit list:
%
% This paper makes the following contributions:
% \begin{itemize}
%   \item \textbf{<Contribution 1>} — one sentence, ending with its headline evidence.
%   ...
% \end{itemize}
% Close with a results preview.

\section{Background and Motivation}
\label{sec:background}
% Anchor on a MOTIVATING MEASUREMENT (a table or figure), not opinion.

\section{<Design / Architecture>}
\label{sec:design}

\section{<The characterization section — the paper's distinctive frame>}
\label{sec:anatomy}
% Write as cohesive original prose. This is usually the most-cited section.

% One section per remaining contribution: mechanism -> why the obvious
% alternative fails -> measured result.

\section{Evaluation}
\label{sec:eval}
% Open with explicit questions: "We answer: Q1 ... Qn".
\subsection{Methodology}
% Testbed table (exact hardware, OS, runtime versions, commit hash, run count,
% warm-up protocol). Named baselines WITH versions. How each metric is computed.
\subsection{Results}
% Medians/error bars over N runs. Every number traces to results/measurements.csv.
\subsection{Ablations}

\section{Related Work}
\label{sec:related}
% Transcribed from the plan doc's verified survey; positioned, not listed.

\section{Limitations and Future Work}
\label{sec:limitations}
% Device variance, thermal, version drift, n, what was NOT measured, author bias.

\section{Conclusion}
\label{sec:conclusion}

\bibliography{refs}

\appendix
\section{Artifact Appendix}
% Repo URL, commit, exact commands to reproduce every table/figure, expected numbers.

\end{document}
